% !TEX program = tectonic
% Slides « Stratégies de trading du Congrès » — vue d'ensemble pour la réunion.
% Format repris de 00_S1S2_donnees/docs/SLIDES_DONNEES_S1S2_V2.tex (mêmes thème/macros).
% Périmètre (décidé avec Alice) : 6 papiers académiques + produits/benchmarks (ETF NANC/GOP,
% Quiver, QuantConnect, Pelosi). Deck STRICTEMENT descriptif/comparatif — « ce qui existe » :
% AUCUNE stratégie perso (pas de ticker V0/nb11, pas de résultats nb06/nb08, pas de Master/Ticker).
% Chiffres = section « Corrections & nuances » du plan (source = PDF primaires /Users/.../Downloads/paper/
% + 00. S3S4 en cours/ETAT_DE_L_ART_STRATEGIES.md). Règles du chef : une idée/slide (titre = message),
% plan annoncé, chronologie, chiffres légendés, choix explicites, public neuf.
\documentclass[11pt,aspectratio=169]{beamer}
\definecolor{helec}{HTML}{3B6EA5}   % couleur maîtresse
\usetheme{Madrid}
\usecolortheme[named=helec]{structure}
\setbeamertemplate{navigation symbols}{}
\setbeamertemplate{blocks}[rounded][shadow=false]
\setbeamerfont{frametitle}{series=\bfseries}
\usepackage[T1]{fontenc}
\usepackage{lmodern}
\usepackage[utf8]{inputenc}
\usepackage[french]{babel}
\usepackage{booktabs}
\usepackage{xcolor}
\usepackage{graphicx}
\usepackage{amssymb}
\usepackage{amsmath}
\usepackage{array}
\usepackage{newunicodechar}
\usepackage{tikz}
\usetikzlibrary{arrows.meta,positioning,calc}
\hypersetup{colorlinks=true, urlcolor=helec, linkcolor=white}

% Caractères Unicode → commandes (fontes T1)
\newunicodechar{→}{\ensuremath{\rightarrow}}
\newunicodechar{—}{\textemdash}
\newunicodechar{–}{\textendash}
\newunicodechar{…}{\ldots}
\newunicodechar{·}{\textperiodcentered}
\newunicodechar{≈}{\ensuremath{\approx}}
\newunicodechar{×}{\ensuremath{\times}}
\newunicodechar{≤}{\ensuremath{\leq}}
\newunicodechar{≥}{\ensuremath{\geq}}
\newunicodechar{−}{\ensuremath{-}}
\newunicodechar{✓}{\checkmark}
\newunicodechar{↔}{\ensuremath{\leftrightarrow}}
\newunicodechar{œ}{\oe}
\newunicodechar{Œ}{\OE}
\newunicodechar{°}{\textdegree}
\newunicodechar{§}{\S}
\newunicodechar{«}{\og}
\newunicodechar{»}{\fg{}}
\newunicodechar{Σ}{\ensuremath{\Sigma}}
\newunicodechar{α}{\ensuremath{\alpha}}
\newunicodechar{β}{\ensuremath{\beta}}
\newunicodechar{∼}{\ensuremath{\sim}}

% Palette
\definecolor{okgreen}{HTML}{2E8B57}
\definecolor{alertred}{HTML}{C0392B}
\definecolor{warnorange}{HTML}{E67E22}
\definecolor{grisC}{RGB}{90,90,90}

\newcommand{\co}[1]{\texttt{#1}}
\tikzset{
  step/.style={rectangle, rounded corners=2pt, draw=#1, line width=0.8pt, fill=#1!8,
               align=center, font=\scriptsize, minimum height=0.72cm},
  bande/.style={rectangle, rounded corners=2pt, draw=#1, line width=0.8pt, fill=#1!8,
                align=center, font=\scriptsize, minimum height=0.62cm},
  arr/.style={-{Latex[length=2mm]}, line width=0.8pt, draw=gray!65},
}

% Tuile chiffre-clé : \tuile{couleur}{gros chiffre}{légende}
\newcommand{\tuile}[3]{%
  \begin{minipage}[t]{0.23\textwidth}\centering
    \colorbox{#1}{\parbox[c][21mm][c]{\dimexpr\linewidth-8pt}{\centering
      {\color{white}\Large\textbf{#2}}\\[0.6mm]
      {\color{white}\footnotesize #3}}}
  \end{minipage}}

\newcommand{\badge}[2]{\colorbox{#1}{\textcolor{white}{\footnotesize\textbf{#2}}}}
\newcommand{\badgel}[2]{\colorbox{#1}{\textcolor{black!80}{\footnotesize\textbf{#2}}}}

% Carte : \carte{couleur}{titre}{contenu}
\newcommand{\carte}[3]{%
  \colorbox{#1!10}{\parbox[t][23mm][t]{\dimexpr\linewidth-8pt}{\vspace{1mm}
    {\color{#1}\footnotesize\textbf{#2}}\\[0.8mm]{\scriptsize #3}}}}

% Encadré CHOIX (un choix méthodologique, en une phrase)
\newcommand{\choixbox}[1]{%
  \colorbox{alertred!8}{\parbox{\dimexpr\linewidth-8pt}{\scriptsize
    \textbf{\textcolor{alertred}{CHOIX}}\; --- #1}}}

% Bandeau RÉPONSE (dernière slide d'une partie)
\newcommand{\repband}[1]{%
  \colorbox{okgreen!12}{\parbox{\dimexpr\textwidth-10pt}{\scriptsize
    \textbf{\textcolor{okgreen}{Réponse}}\; --- #1}}}

% Bandeau TENSION (question ouverte)
\newcommand{\tension}[1]{%
  \colorbox{warnorange!12}{\parbox{\dimexpr\textwidth-10pt}{\scriptsize
    \textbf{\textcolor{warnorange}{Ouvert}}\; --- #1}}}

% Ligne d'un sommaire de partie : \etape{n}{titre}{glose}
\newcommand{\etape}[3]{%
  {\color{helec}\bfseries #1}\; & \textbf{#2} & {\scriptsize\color{black!60} #3}\\[0.6mm]}

\title[Copier le Congrès]{Stratégies de trading du Congrès américain}
\subtitle{Ce que disent les recherches \& les produits existants · vue d'ensemble pour la réunion}
\author[A.~Lemaire]{Alice Lemaire}
\institute[Ramify]{Document de recherche --- livrable Ramify}
\date{Juillet 2026}

\begin{document}

% ── Titre ──
\begin{frame}[plain]
  \titlepage
\end{frame}

% ══════════════════ OUVERTURE ══════════════════

% ── Le plan ──
\begin{frame}{Le plan --- cinq questions, dans l'ordre de vie d'un trade}
  \centering
  \vspace{1mm}
  {\footnotesize la même donnée --- un trade d'élu déclaré publiquement --- lue sous cinq angles successifs}\\[4mm]
  {\small\setlength{\extrarowheight}{2.2pt}
  \begin{tabular}{@{}p{0.8cm}>{\raggedright\arraybackslash}p{5.0cm}>{\raggedright\arraybackslash}p{7.0cm}@{}}
    \etape{I}{Le signal, c'est la divulgation}{Lazzaretto, Pyun --- et non la transaction}
    \etape{II}{D'où vient l'avantage}{Wei \& Zhou (chefs de parti), Li (jalons de loi)}
    \etape{III}{Deux angles alternatifs}{Molk \& Partnoy (paris baissiers), Roodman (ML sur graphe)}
    \etape{IV}{Ce qui existe déjà}{ETF NANC/GOP, Quiver, QuantConnect, Pelosi}
    \etape{V}{Synthèse \& pistes}{ce qui est robuste, ce qui reste ouvert}
  \end{tabular}}
\end{frame}

% ── Glossaire ──
\begin{frame}{Onze mots à connaître avant de commencer}
  \centering
  {\footnotesize le seul chiffre à retenir ici : \textbf{$\leq$ 45 jours} entre le trade et sa publication}\\[2.5mm]
  \begin{columns}[T]
    \begin{column}{0.49\textwidth}
      {\scriptsize
      \badge{helec}{STOCK Act (2012)} loi obligeant tout élu à déclarer ses trades\\[1.4mm]
      \badge{helec}{PTR} le formulaire de déclaration (\emph{Periodic Transaction Report})\\[1.4mm]
      \badge{grisC}{Date de transaction} le jour où l'élu achète/vend (privé)\\[1.4mm]
      \badge{helec}{Date de divulgation} le jour où le PTR devient public ($\leq$ 45 j après)\\[1.4mm]
      \badge{okgreen}{Alpha ($\alpha$)} le rendement au-delà de ce qu'explique le marché\\[1.4mm]
      \badge{okgreen}{BHAR} sur-rendement cumulé d'un titre, acheté et conservé
      }
    \end{column}
    \begin{column}{0.49\textwidth}
      {\scriptsize
      \badge{okgreen}{CAAR / CAR} sur-rendement moyen cumulé autour d'un événement\\[1.4mm]
      \badge{grisC}{Value-weighted} pondéré par la taille : les gros trades pèsent plus\\[1.4mm]
      \badge{grisC}{Long-short} acheter les uns, vendre à découvert les autres\\[1.4mm]
      \badge{warnorange}{Stealth trading} fractionner un gros ordre pour ne pas se faire repérer\\[1.4mm]
      \badge{warnorange}{Enrollment} texte final voté par les 2 chambres --- un « jalon » de loi ; un \emph{enterprising trade} est un trade juste avant un tel jalon
      }
    \end{column}
  \end{columns}
\end{frame}

% ══════════════════ PARTIE I ══════════════════
\section{Partie I · Le signal, c'est la divulgation}

\begin{frame}[plain]
  \centering
  \vspace{12mm}
  {\color{helec!85}\large\textbf{PARTIE I}}\\[3mm]
  {\color{black!88}\fontsize{30}{34}\selectfont\textbf{Le signal, c'est la divulgation}}\\[3mm]
  {\color{helec!45}\rule{0.30\textwidth}{1.2pt}}\\[3mm]
  {\color{black!55}\large Pas la transaction --- et c'est ce qui le rend copiable}
\end{frame}

% ── S4 · Deux dates ──
\begin{frame}{Le STOCK Act crée deux dates --- et l'anomalie n'est qu'à la seconde}
  \centering
  \vspace{0.5mm}
  {\footnotesize\textit{Question de la Partie~I --- où est l'événement qu'on peut vraiment exploiter ?}}\\[5mm]
  \begin{tikzpicture}[node distance=32mm]
    \node[step=grisC, text width=3.7cm] (t)
      {\textbf{Date de transaction}\\[0.4mm]{\scriptsize l'élu achète / vend\\ privé, invisible}};
    \node[step=helec, text width=3.7cm, right=of t] (d)
      {\textbf{Date de divulgation}\\[0.4mm]{\scriptsize le PTR devient public\\ copiable par tous}};
    \draw[arr] (t) -- node[above,font=\scriptsize]{$\leq$ 45 j}
                       node[below,font=\scriptsize]{médiane $\sim$27 j} (d);
  \end{tikzpicture}\\[5mm]
  {\footnotesize En agrégé, le sur-rendement exploitable après la \textbf{seule transaction} est $\approx$~0
   (Lazzaretto) :\\ l'information ne devient actionnable qu'à la \textbf{divulgation}.}\\[3.5mm]
  \choixbox{on date tout à la \textbf{divulgation}, jamais à la transaction --- sinon on utiliserait
   une information pas encore publique (\emph{look-ahead}).}
\end{frame}

% ── S5 · Lazzaretto ──
\begin{frame}{À la divulgation, le prix bouge --- puis se rétracte en 60 jours}
  \centering
  {\footnotesize\textbf{Lazzaretto (2024)} --- \emph{Mind the Trade}, Sénat, 2018--2022}\\[4mm]
  \tuile{helec}{+90 bps}{achats : le mois\\qui suit la divulgation}\hfill
  \tuile{alertred}{$-$60 bps}{ventes : même\\fenêtre, le prix baisse}\hfill
  \tuile{warnorange}{$\sim$60 j}{puis tout se rétracte :\\réversion en $\sim$60 j de bourse}\\[5mm]
  {\footnotesize Le portefeuille \textbf{long-short} (achats $-$ ventes), en calendar-time, dégage
   $>$~2~\%/mois brut ---\\ mais ce rendement s'\textbf{inverse} : c'est une fenêtre courte, pas un alpha qu'on conserve.}
\end{frame}

% ── S6 · Pyun (clôture I) ──
\begin{frame}{Après divulgation, il reste un alpha --- plus mince, côté grandes caps et Chambre}
  \centering
  {\footnotesize\textbf{Pyun (2026)} --- portefeuille d'achats divulgués, pondéré-valeur, 2014--2023}\\[4mm]
  \tuile{helec}{$\approx$0,3 \%}{alpha mensuel, hors\\coûts (FF / q-factor)}\hfill
  \tuile{okgreen}{+3 à 4 \%}{excès annuel du\\portefeuille vs S\&P}\hfill
  \tuile{grisC}{16,5 vs 14,4}{rendement à 240 j (\%) :\\dater sur la transaction $>$ report}\\[4mm]
  {\footnotesize Concentré là où l'information se diffuse mal : \textbf{large-caps}, \textbf{Représentants}
   (77,7~\% des transactions), \textbf{post-COVID}.}\\[3.5mm]
  \repband{copiable, oui --- mais \textbf{court}, \textbf{réversif} et net de peu : le Sharpe relatif
   (achats VW $-$ S\&P) n'est que \textbf{0,66}, avant coûts de transaction.}
\end{frame}

% ══════════════════ PARTIE II ══════════════════
\section{Partie II · D'où vient l'avantage}

\begin{frame}[plain]
  \centering
  \vspace{12mm}
  {\color{helec!85}\large\textbf{PARTIE II}}\\[3mm]
  {\color{black!88}\fontsize{30}{34}\selectfont\textbf{D'où vient l'avantage}}\\[3mm]
  {\color{helec!45}\rule{0.30\textwidth}{1.2pt}}\\[3mm]
  {\color{black!55}\large L'edge n'est pas partout : il se concentre chez certains élus, à certains moments}
\end{frame}

% ── S7 · Wei & Zhou ──
\begin{frame}{L'avantage n'est pas partout --- il est chez les chefs de parti}
  \centering
  {\footnotesize\textit{Question de la Partie~II --- qui a un vrai edge, et pourquoi ?}\quad
   \textbf{Wei \& Zhou (2025)}, NBER, 47 leaders, 1995--2021}\\[4mm]
  \tuile{helec}{+47 pp}{« jusqu'à », par an vs\\pairs appariés (20 leaders)}\hfill
  \tuile{okgreen}{+29 / +24}{pp à 6 mois :\\achats / ventes des chefs}\hfill
  \tuile{grisC}{+7,5 bps}{par jour, sur-rendement\\après l'ascension}\\[5mm]
  \choixbox{« leadership » = 12 postes de direction (Speaker, floor/whip, conference) ---
   \emph{pas} la simple appartenance à une commission.}
\end{frame}

% ── S8 · Li (clôture II) ──
\begin{frame}{L'autre source d'edge : trader juste avant un jalon législatif}
  \centering
  {\footnotesize\textbf{Li et al.\ (2026)} --- \emph{Inside the Beltway}, Sénat, 2014--2021 ;
   les élus tradent 1--30 j \textbf{avant} ces jalons}\\[3.5mm]
  \begin{tikzpicture}[node distance=4mm]
    \node[bande=grisC, text width=2.3cm] (a) {introduction};
    \node[bande=grisC, text width=2.3cm, right=of a] (b) {engrossment};
    \node[bande=helec, text width=2.3cm, right=of b] (c) {\textbf{enrollment}};
    \node[bande=grisC, text width=2.3cm, right=of c] (e) {enactment};
    \draw[arr] (a) -- (b); \draw[arr] (b) -- (c); \draw[arr] (c) -- (e);
  \end{tikzpicture}\\[4mm]
  \tuile{helec}{+1,16 \%}{coefficient au stade\\« enrollment » (coef.\ +1,159)}\hfill
  \tuile{okgreen}{+6,21 \%}{si l'élu est très efficace\\législativement (High-LES)}\hfill
  \tuile{warnorange}{$\sim$22 \%}{des trades en séquences\\de 2 à 5 (stealth)}\\[3.5mm]
  \repband{deux edges \textbf{réels mais conditionnels} : peu d'élus concernés, et il faut des
   données de contexte externes (leadership, calendrier des lois).}
\end{frame}

% ══════════════════ PARTIE III ══════════════════
\section{Partie III · Deux angles alternatifs}

\begin{frame}[plain]
  \centering
  \vspace{12mm}
  {\color{helec!85}\large\textbf{PARTIE III}}\\[3mm]
  {\color{black!88}\fontsize{30}{34}\selectfont\textbf{Deux angles alternatifs}}\\[3mm]
  {\color{helec!45}\rule{0.30\textwidth}{1.2pt}}\\[3mm]
  {\color{black!55}\large Et si on cherchait au mauvais endroit --- ou avec le mauvais outil ?}
\end{frame}

% ── S9 · Molk & Partnoy ──
\begin{frame}{Le vrai edge est peut-être baissier --- et c'est notre angle mort}
  \centering
  {\footnotesize\textit{Question de la Partie~III --- et si on cherchait au mauvais endroit ?}\quad
   \textbf{Molk \& Partnoy (2025)}, $\sim$30 élus, 2012--2023}\\[4mm]
  \tuile{helec}{587}{trades « négatifs »\\(puts, ventes à découvert)}\hfill
  \tuile{alertred}{$-$0,55 \%}{CAR[$-$1,1] à la divulgation :\\le pari baissier gagne}\hfill
  \tuile{grisC}{$\approx$0}{les positions longues,\\elles, ne battent rien}\\[5mm]
  {\footnotesize Volume concerné : \textbf{$\sim$15,6~M\$} (options : 414 trades / $\sim$9,2~M\$ ;
   ventes à découvert : 116 / 2,08~M\$).}\\[3mm]
  \choixbox{ces paris passent par \textbf{options} et \textbf{ventes à découvert} : hors de notre
   périmètre actions --- angle mort assumé.}
\end{frame}

% ── S10 · Roodman (clôture III) ──
\begin{frame}{Un graphe temporel de ML ne prédit pas mieux que pile ou face}
  \centering
  {\footnotesize\textbf{Roodman et al.\ (2026)} --- GAP-TGN, graphe multimodal, 32\,401 trades, 2013--2025}\\[4mm]
  \tuile{alertred}{0,52}{AUROC à 24 mois :\\0,50 = pur hasard}\hfill
  \tuile{grisC}{0,44 vs 0,29}{F1 vs XGBoost --- mais les\\baselines s'effondrent}\hfill
  \tuile{helec}{4 modalités}{trades + lobbying\\+ votes + géographie}\\[5mm]
  {\footnotesize Le modèle \textbf{ne discrimine pas} les trades gagnants des perdants : un AUROC de 0,52
   est indiscernable d'un tirage à pile ou face.}\\[3.5mm]
  \repband{ni le short ni le ML ne renversent le verdict : l'essentiel des trades n'a \textbf{pas} de
   signal exploitable --- ces deux angles \textbf{délimitent} nos angles morts, sans les combler.}
\end{frame}

% ══════════════════ PARTIE IV ══════════════════
\section{Partie IV · Ce qui existe déjà}

\begin{frame}[plain]
  \centering
  \vspace{12mm}
  {\color{helec!85}\large\textbf{PARTIE IV}}\\[3mm]
  {\color{black!88}\fontsize{30}{34}\selectfont\textbf{Ce qui existe déjà}}\\[3mm]
  {\color{helec!45}\rule{0.30\textwidth}{1.2pt}}\\[3mm]
  {\color{black!55}\large Les produits déjà investissables --- et ce que valent leurs chiffres}
\end{frame}

% ── S11 · ETF NANC/GOP ──
\begin{frame}{Les ETF encadrent le S\&P --- NANC devant, GOP derrière}
  \centering
  {\footnotesize\textit{Question de la Partie~IV --- qu'obtient-on vraiment en copiant, net de frais ?}}\\[3.5mm]
  \tuile{okgreen}{+18,66 \%}{NANC · 1 an\\(S\&P TR : +17,88 \%)}\hfill
  \tuile{helec}{+23,63 \%}{NANC · création\\(S\&P : +20,96 \%)}\hfill
  \tuile{warnorange}{+17,16 \%}{GOP · 1 an\\(S\&P TR : +17,88 \%)}\hfill
  \tuile{alertred}{+14,75 \%}{GOP · création\\(S\&P : +20,96 \%)}\\[2.5mm]
  {\scriptsize\color{black!60} Rendements \emph{audités} (prospectus), périodes closes au \textbf{31/12/2025} ·
   NANC réplique surtout les Démocrates (forte tech, top-10 $\approx$ 46~\%)}\\[2.5mm]
  {\footnotesize Mécanique commune : lookback $\sim$3 ans, prix milieu de fourchette, 100--200 lignes,
   on \textbf{nette} achats $-$ ventes, on surpondère gros / récurrents / multi-élus.}\\[2.5mm]
  \choixbox{benchmark = \textbf{S\&P 500 Total Return}, et on légende toujours la date d'arrêté
   (à mi-2026, le 1 an de NANC repasse \emph{sous} le S\&P).}
\end{frame}

% ── S12 · Backtests publics (clôture IV) ──
\begin{frame}{Les backtests publics brillent parce qu'ils ignorent les coûts (et les options)}
  \centering
  \vspace{0.5mm}
  \tuile{helec}{37 \%}{Quiver : CAGR d'un backtest\\« brut » ($\beta$ 1,14, \textbf{0 coût})}\hfill
  \tuile{warnorange}{$\sim$14 \%/an}{Pelosi, cumulé « 40 ans »\\(S\&P $\sim$9 \%) --- via options}\hfill
  \tuile{grisC}{$\Sigma$B$-\Sigma$V}{QuantConnect : un \textbf{signal}\\(net acheteur → long)}\\[3mm]
  {\footnotesize Le miroir \textbf{actions seul} de Pelosi (Quiver) ne donne qu'un Sharpe de \textbf{0,72} :
   l'essentiel du rendement vient du \textbf{levier optionnel}, pas des actions.\\
   QuantConnect ne publie \emph{aucune} performance --- c'est un modèle pédagogique.}\\[3mm]
  \choixbox{ces chiffres sont \textbf{bruts} : à répliquer nets de coûts et de turnover, avec un test
   de significativité, avant d'y croire.}\\[2mm]
  \repband{en réel, net de frais, copier le Congrès rapporte $\approx$ le marché --- NANC lui-même est
   surtout un \textbf{beta} tech (top-10 $\approx$ 46~\%).}
\end{frame}

% ══════════════════ PARTIE V ══════════════════
\section{Partie V · Synthèse \& pistes}

\begin{frame}[plain]
  \centering
  \vspace{12mm}
  {\color{helec!85}\large\textbf{PARTIE V}}\\[3mm]
  {\color{black!88}\fontsize{30}{34}\selectfont\textbf{Synthèse \& pistes}}\\[3mm]
  {\color{helec!45}\rule{0.30\textwidth}{1.2pt}}\\[3mm]
  {\color{black!55}\large Ce qui est robuste, ce qui reste ouvert, et par où continuer}
\end{frame}

% ── S13 · Tableau récapitulatif ──
\begin{frame}{Six approches (plus les produits), une grille de décision}
  \centering
  {\scriptsize\setlength{\tabcolsep}{3.5pt}\setlength{\extrarowheight}{1pt}
  \begin{tabular}{@{}>{\raggedright\arraybackslash}p{1.7cm}>{\raggedright\arraybackslash}p{2.9cm}>{\raggedright\arraybackslash}p{1.5cm}>{\raggedright\arraybackslash}p{1.4cm}>{\raggedright\arraybackslash}p{3.0cm}>{\raggedright\arraybackslash}p{1.9cm}@{}}
    \toprule
    \textbf{Approche} & \textbf{Idée} & \textbf{Daté sur} & \textbf{Horizon} & \textbf{Chiffre clé} & \textbf{Réplicable} \\
    \midrule
    Lazzaretto & La divulgation bouge les prix, pas la transaction & Divulgation & court ($<$60 j) & +90 bps/mo (achats), réverse & Oui, horizon court \\
    Pyun & Portefeuille d'achats VW divulgués & Transaction \& report & 1--12 mois & $\alpha$ $\sim$0,3 \%/mo ; Sharpe 0,66\textsuperscript{$\ast$} & Oui (le plus direct) \\
    Wei \& Zhou & Suivre les chefs de parti & Trade + ascension & 6--12 mois & jusqu'à +47 pp/an & Partiel (canaux ext.) \\
    Li et al. & Trader avant les jalons de loi & Trade vs jalon & $\leq$ 1 an & +1,16 \% ; +6,21 \% (High-LES) & Non / partiel \\
    Molk \& Partnoy & Copier les paris baissiers & Divulgation & 10--15 j & CAR $-$0,55 \% ; longs $\approx$0 & Non (options/short) \\
    Roodman & ML sur graphe multimodal & Divulgation & 18--24 mois & AUROC 0,52 $\approx$ hasard & Confirme : pas d'alpha \\
    \midrule
    ETF NANC/GOP & Produit qui réplique & Divulgation & continu & NANC +23,63 \%/an (création) & Déjà investissable \\
    Quiver / QC / Pelosi & Trackers \& backtests médiatiques & mixte & divers & Quiver 37 \% (0 coût) ; Pelosi options & À répliquer proprement \\
    \bottomrule
  \end{tabular}}\\[1mm]
  {\scriptsize\color{black!60} \textsuperscript{$\ast$} Sharpe relatif (achats VW $-$ S\&P), avant coûts.}
\end{frame}

% ── S14 · Ce qui est robuste ──
\begin{frame}{Ce qui ressort de robuste --- le consensus des recherches}
  \centering
  \vspace{1mm}
  \begin{columns}[T]
    \begin{column}{0.32\textwidth}
      \carte{helec}{1 · Achats $>$ ventes}{le signal acheteur domine partout ; les ventes portent peu (sauf paris baissiers)}
    \end{column}
    \begin{column}{0.32\textwidth}
      \carte{helec}{2 · Grandes caps}{l'alpha survit surtout sur les large-caps (Pyun)}
    \end{column}
    \begin{column}{0.32\textwidth}
      \carte{helec}{3 · Netter A $-$ V}{achats moins ventes, sur un lookback $\sim$3 ans (Pyun, ETF)}
    \end{column}
  \end{columns}\par\vspace{3mm}
  \begin{columns}[T]
    \begin{column}{0.32\textwidth}
      \carte{helec}{4 · Pondérer-valeur}{prix milieu de fourchette ; les gros trades pèsent plus}
    \end{column}
    \begin{column}{0.32\textwidth}
      \carte{helec}{5 · Signaux de qualité}{comité (Pyun/Li), leadership (Wei\,\&\,Zhou), récurrence / multi-élus (ETF)}
    \end{column}
    \begin{column}{0.32\textwidth}
      \carte{helec}{6 · Dater la divulgation}{c'est là que le prix réagit, pas à la transaction}
    \end{column}
  \end{columns}
\end{frame}

% ── S15 · Ce qui reste ouvert ──
\begin{frame}{Ce qui reste ouvert --- à traiter avec prudence}
  {\small
  \begin{itemize}\setlength{\itemsep}{1.3mm}
    \item[\textcolor{warnorange}{$\bullet$}] \textbf{Réversion vs drift} : Lazzaretto voit le signal se rétracter en $<$60 j ; Pyun et le graphe ML supposent un drift à 18--24 mois --- l'horizon reste à trancher.
    \item[\textcolor{warnorange}{$\bullet$}] \textbf{Coûts non nets} : presque tous les chiffres sont \emph{bruts} (0 coût de transaction, 0 turnover).
    \item[\textcolor{warnorange}{$\bullet$}] \textbf{Échantillons courts} : 2018--2022 pour les papiers ; les ETF n'existent que depuis 2023.
    \item[\textcolor{warnorange}{$\bullet$}] \textbf{Look-ahead} : dater sur la transaction utilise une information pas encore publique.
    \item[\textcolor{warnorange}{$\bullet$}] \textbf{Corroboration = signal ou herding ?} plusieurs acheteurs = plus fiable (ETF), ou au contraire moins performant (Herron) ?
    \item[\textcolor{warnorange}{$\bullet$}] \textbf{Consensus historique mitigé} : Eggers--Hainmueller, Belmont ne trouvent pas d'alpha global.
  \end{itemize}}
  \vspace{1mm}
  \tension{ces poches de rendement sont des \textbf{exceptions conditionnelles} (certains élus, certains moments), \emph{pas} une réfutation du « pas d'alpha » général.}
\end{frame}

% ── S16 · Pistes ──
\begin{frame}{Pistes pour aller plus loin}
  \centering
  \vspace{1mm}
  \begin{columns}[T]
    \begin{column}{0.32\textwidth}
      \carte{helec}{a · Event-study double-date}{mesurer autour de la transaction \emph{et} de la divulgation --- pour trancher réversion vs drift}
    \end{column}
    \begin{column}{0.32\textwidth}
      \carte{okgreen}{b · Répliquer NANC/GOP proprement}{avec coûts, puis régresser vs SPY + panier sectoriel GICS → séparer \textbf{alpha} et \textbf{beta}}
    \end{column}
    \begin{column}{0.32\textwidth}
      \carte{grisC}{c · Fixer les benchmarks}{cibles claires : NANC ($+$23,63 \%/an depuis création) et le S\&P 500 (SPY)}
    \end{column}
  \end{columns}\par\vspace{5mm}
  \repband{l'objectif n'est pas de « battre » dans un backtest brut, mais de séparer le \textbf{vrai alpha}
   du \textbf{beta déguisé} --- net de coûts.}
\end{frame}

\end{document}
